\documentclass[12pt,a4paper]{article}
\usepackage[utf8]{inputenc}
\usepackage[T1]{fontenc}
\usepackage[turkish]{babel}
\usepackage{lmodern}
\usepackage{microtype}
\usepackage{amsmath, amssymb}
\usepackage{geometry}
\usepackage{graphicx}
\usepackage{setspace}
\usepackage{cite}
\usepackage{titlesec}
\usepackage[colorlinks=true, linkcolor=blue, urlcolor=blue, citecolor=black]{hyperref}

\geometry{margin=2.5cm}
\setstretch{1.3} % Okunabilirliği artıran satır aralığı
\titleformat{\section}{\bfseries\Large}{\thesection.}{1em}{}
\titlespacing*{\section}{0pt}{1.5ex plus 0.5ex minus 0.2ex}{0.8ex plus 0.2ex}
\titleformat{\subsection}{\bfseries\normalsize}{\thesubsection.}{1em}{}
\titlespacing*{\subsection}{0pt}{1.2ex plus 0.4ex minus 0.2ex}{0.5ex}

\title{\textbf{Duran Sicimler Teoremi: Görelilik, Faz Dönüşümü ve Bilginin Kozmolojik Döngüsü}\\
\large (Görelilik ve Duran Sicimler Arasındaki Köprü – Genişletilmiş Versiyon)}
\author{Fahri Meral}
\date{2025}

\begin{document}

\maketitle

\begin{abstract}
Bu çalışma, evreni iki temel varoluş hâli üzerinden yeniden yorumlamayı amaçlamaktadır: \textit{titreşen enerji fazı} ve \textit{duran bilgi fazı}. 
Einstein’ın genel görelilik kuramı, güçlü kütleçekim alanlarında zamanın yavaşladığını ve olay ufkunda durduğunu öngörür. 
Sicim kuramı ise evrenin temelinde titreşen enerji sicimlerinin bulunduğunu ileri sürer. 
Bu iki bakış açısını birleştiren \textbf{Duran Sicimler Teoremi}, zamanın akışını sicimlerin titreşim frekansına ($f$) bağlar. 

Bu çalışmada, düzenin ölçüsü olarak ``Mükemmellik Potansiyeli'' ($\Phi$) adında yeni bir fiziksel büyüklük önerilmektedir. Bu potansiyel, entropinin ($S$) tersi olarak tanımlanır: $\Phi = 1/S(f)$. Kara delikler, sicim titreşiminin durduğu ($f \to 0$) ve entropinin sıfıra yaklaştığı ($S \to 0$) kozmik mekanizmalardır; bu durumda Mükemmellik Potansiyeli sonsuza ıraksar ($\Phi \to \infty$). Bu faz dönüşümü, enerjinin yok olmadığını, yalnızca mutlak düzenin hüküm sürdüğü durağan bilgiye dönüştüğünü gösterir. Böylece evren, bilgi ve enerji arasında sürekli bir döngü içinde var olur.
\end{abstract}

\section{Giriş}
Albert Einstein, genel görelilik kuramıyla zamanın mutlak değil, kütleçekim alanı tarafından eğrilen bir boyut olduğunu göstermiştir. 
Benzer biçimde, sicim teorisi de evrendeki tüm parçacıkların farklı frekanslarda titreşen sicimlerden oluştuğunu öne sürer. 
Her iki yaklaşımın da temelinde \textit{zaman–enerji ilişkisi} yatar.

Bu makale, bu iki çerçeveyi birleştirerek şu temel önermeyi ortaya koyar:
\begin{quote}
``Zamanın yavaşlaması, sicimlerin titreşim frekansının azalmasının doğrudan sonucudur. 
Zamanın durması, sicimlerin titreşimini yitirip mükemmel düzenin ve saf bilginin hüküm sürdüğü durağan hâle geçmesidir.''
\end{quote}

Bu varsayıma göre kara delikler, evrende bir \textit{faz dönüşümü} gerçekleştirir: 
Titreşen enerji formu, durağan bilgi formuna geçer. 
Bu bilgi, enerjinin varoluşsal izidir; yani evrenin ``bilgisel DNA''sıdır.

\section{Görelilikte Zaman Genişlemesi}

Genel görelilikte zaman, kütleçekim potansiyeliyle doğrudan ilişkilidir. 
Bir gözlemci kara deliğe yaklaştıkça zamanın akışı yavaşlar ve olay ufkunda tamamen durur.

Zaman genişlemesi denklemi şöyledir:
\begin{equation}
t' = t \sqrt{1 - \frac{2GM}{r c^2}}
\end{equation}

Burada:
\begin{itemize}
\item $t'$ : yerel gözlemcinin ölçtüğü zaman,
\item $t$ : uzak gözlemcinin ölçtüğü zaman,
\item $G$ : evrensel çekim sabiti,
\item $M$ : kara deliğin kütlesi,
\item $r$ : gözlemcinin kara deliğin merkezine uzaklığı,
\item $c$ : ışık hızı.
\end{itemize}

$r \to \frac{2GM}{c^2}$ olduğunda, kökün içi sıfıra yaklaşır ve $t' \to 0$ olur. 
Bu, olay ufkunda zamanın durması anlamına gelir.

\subsection{Olay Ufku ve Zamanın Yavaşlaması}
Olay ufku, yalnızca ışığın kaçamadığı bir sınır değil, zamanın anlamını kaybettiği bir noktadır. 
Bu sınırın ötesinde klasik zaman kavramı çöker; bu da sicimsel titreşimin yavaşlamasıyla doğrudan ilişkilendirilebilir.

\section{Sicimlerin Frekansı ile Zamanın Bağıntısı}
Sicim teorisinde enerji, sicimin frekansıyla doğru orantılıdır:
\begin{equation}
E = h f
\end{equation}
ve Einstein kütle-enerji denkliğiyle birlikte:
\begin{equation}
E = mc^2 \Rightarrow m = \frac{h f}{c^2}
\end{equation}
Buradan sicimin titreşim frekansı azaldıkça enerjisinin de azaldığı görülür.

Kütleçekim altında frekans, zaman genişlemesine benzer biçimde değişir:
\begin{equation}
f' = f_0 \sqrt{1 - \frac{2GM}{r c^2}}
\end{equation}
Olay ufkunda $f' \to 0$ olur; yani sicim titreşimini yitirir.

\subsection{Frekansın Kütleçekime Göre Değişimi}
Bu ifade, zamanla frekans arasında doğrudan bir bağ kurulabileceğini gösterir:
\begin{equation}
t \propto f
\end{equation}
Dolayısıyla, frekansın sıfıra yaklaşması zamanın da durmasına yol açar.

\section{Birleştirilmiş Denklem ve Zamanın Durağanlığı}
Zaman ve sicim frekansını birleştiren genel formül:
\begin{equation}
t' = k_t \cdot f_0 \sqrt{1 - \frac{2GM}{r c^2}}
\end{equation}
Burada $k_t$ sabiti, zamanın algısal ölçeğini temsil eder. 
$r \to \frac{2GM}{c^2}$ olduğunda $t' \to 0$ ve $f' \to 0$ olur; 
yani sicim titreşimi durur, zaman da anlamını yitirir.

\section{Evrenin İki Fazı: Enerji ve Bilgi}

\subsection{Titreşen Enerji Fazı}
Bu faz, gözlemlenebilir evrendir: sicimler titreşir ($f > 0$), enerji ve madde etkileşir, entropi artar. 
Zamanın oku ileriye doğrudur. 
Bu fazda düzen sürekli bozulur ve evren giderek daha yüksek entropiye ulaşır.

\subsection{Duran Bilgi Fazı (Mükemmellik Boyutu)}
Bu fazda sicimlerin titreşimi durur ($f = 0$). 
Enerji, zaman ve entropi kavramları anlamını yitirir. 
Bu, maddenin dönüşeceği ``mükemmel hâl''dir. 
Bilgi mutlak düzen hâline gelir — sistemin tek bir mikro-durumu vardır. 

\subsection{Mükemmellik Potansiyeli ($\Phi$) ve Matematiksel Model}
Evrendeki düzeni ve kaosu ölçmek için ``Mükemmellik Potansiyeli'' ($\Phi$) olarak adlandırdığımız yeni bir fiziksel büyüklük tanımlayabiliriz. Bu potansiyel, sistemin sicim frekansına bağlı entropi fonksiyonunun ($S(f)$) tersidir:
\begin{equation}
\Phi = \frac{1}{S(f)}
\end{equation}
Bu tanıma göre:
\begin{itemize}
    \item Titreşim arttıkça ($f > 0$), entropi artar ($S$ artar) ve Mükemmellik Potansiyeli azalır ($\Phi$ azalır).
    \item Titreşim azaldıkça ($f \to 0$), entropi azalır ($S \to 0$) ve Mükemmellik Potansiyeli artar ($\Phi \to \infty$).
\end{itemize}
Kara delikler bu geçişi sağlayan faz dönüştürücülerdir. Olay ufkunda zamanın durması ($t' \to 0$) ve frekansın sıfırlanmasıyla ($f' \to 0$):
\begin{equation}
\lim_{f \to 0} S(f) = 0 \implies \Phi \to \infty
\end{equation}
Bu, entropinin sıfır ve potansiyelin sonsuz olduğu mutlak düzen hâlidir.

\section{Felsefi Yorum: Zamanın Ontolojik Sessizliği}
Eğer zaman, sicimlerin titreşimiyle ortaya çıkan bir olgusa, 
o hâlde zamanın durması, varoluşun sessizliğe dönüşmesi demektir. 
Bu noktada, evren ``hareketsiz mükemmellik'' düzeyine ulaşır. 

Bu yaklaşım, Wheeler’ın ``It from Bit'' önermesiyle de örtüşür: 
varlık (``It''), bilgiden (``Bit'') doğar. 
Bizim modelimizde bu iki hâl arasında bir döngü vardır:
\[
\text{Enerji (Titreşim)} \rightleftharpoons \text{Bilgi (Durağanlık)}
\]

\section{Kozmolojik Döngü ve Sonuç}
Evrenin döngüsel bir yapıya sahip olduğu varsayılır:
\begin{itemize}
\item \textbf{Başlangıç (Büyük Patlama):} Durağan bilgi fazının ($\Phi \to \infty$) enerjiye dönüşümü ($f > 0$).
\item \textbf{Evrim:} Sicimlerin titreştiği, zamanın aktığı, entropinin arttığı ($S > 0$) dönem.
\item \textbf{Dönüşüm:} Kara delikler, maddeyi yeniden durağan bilgi fazına ($f \to 0, S \to 0, \Phi \to \infty$) çevirir.
\item \textbf{Yeniden Doğuş:} Bilgi tekrar enerjiye dönüşerek yeni bir evreni başlatabilir.
\end{itemize}

Bu çalışma, Einstein’ın zaman genişlemesi denklemine sicimsel bir açıklama getirerek ve Mükemmellik Potansiyeli kavramını tanıtarak, kara delikleri birer faz dönüştürücü olarak tanımlar. 
Zamanın yavaşlaması ve durması, yalnızca göreliliksel bir etki değil; sicimlerin enerjiden ($\Phi \to 0$) bilgiye ($\Phi \to \infty$) geçişinin göstergesidir.

\begin{quote}
\textit{``Evrenin titreşimi yavaşladığında zaman da yavaşlar;\\
zaman durduğunda ise mükemmellik başlar.''}\\
— Fahri Meral, Duran Sicimler Teoremi (2025)
\end{quote}

\section*{Kaynakça}
\begin{enumerate}
\item Meral, F. (2025). \textit{Duran Sicimler Teoremi: Zamanın ve Entropinin Yokluğunda Mükemmel Düzenin Yeniden Doğuşu.}
\item Einstein, A. (1916). \textit{Die Grundlage der allgemeinen Relativitätstheorie}. Annalen der Physik, 49, 769–822.
\item Schwarzschild, K. (1916). \textit{\"Uber das Gravitationsfeld eines Massenpunktes nach der Einsteinschen Theorie}. Sitzungsberichte der Königlich Preussischen Akademie der Wissenschaften.
\item Hawking, S. W. (1974). \textit{Black hole explosions?} Nature, 248(5443), 30–31.
\item Hawking, S. W., \& Ellis, G. F. R. (1973). \textit{The Large Scale Structure of Space-Time}. Cambridge University Press.
\item Bekenstein, J. D. (1973). \textit{Black holes and entropy}. Physical Review D, 7(8), 2333–2346.
\item Maldacena, J. (1998). \textit{The Large N Limit of Superconformal Field Theories and Supergravity}. Advances in Theoretical and Mathematical Physics, 2, 231–252.
\item Maldacena, J., \& Susskind, L. (2013). \textit{Cool horizons for entangled black holes}. Fortschritte der Physik, 61(9), 781–811. (ER=EPR önerisi)
\item Penrose, R. (2010). \textit{Cycles of Time: An Extraordinary New View of the Universe}. Bodley Head.
\item Polchinski, J. (1998). \textit{String Theory Vol. 1 \& 2}. Cambridge University Press.
\item Green, M. B., Schwarz, J. H., \& Witten, E. (1987). \textit{Superstring Theory}. Cambridge University Press.
\item Kaku, M. (1999). \textit{Introduction to Superstrings and M-Theory}. Springer.
\item Greene, B. (1999). \textit{The Elegant Universe: Superstrings, Hidden Dimensions, and the Quest for the Ultimate Theory}. W. W. Norton \& Company.
\item Susskind, L. (2008). \textit{The Black Hole War: My Battle with Stephen Hawking to Make the World Safe for Quantum Mechanics}. Little, Brown and Company.
\item Wheeler, J. A. (1990). \textit{Information, Physics, Quantum: The Search for Links}. In W. H. Zurek (Ed.), \textit{Complexity, Entropy, and the Physics of Information} (pp. 1–28). Addison-Wesley.
\item Page, D. N. (1993). \textit{Information in black hole radiation}. Physical Review Letters, 71(23), 3743–3746.
\item 't Hooft, G. (1993). \textit{Dimensional reduction in quantum gravity}. arXiv:gr-qc/9310026. (Holographic principle)
\item Srednicki, M. (1993). \textit{Entropy and area}. Physical Review Letters, 71(5), 666–669.
\item Bousso, R. (2002). \textit{The holographic principle}. Reviews of Modern Physics, 74(3), 825–874.
\item Preskill, J. (1992). \textit{Do black holes destroy information?} Proceedings of the International Symposium on Black Holes, Membranes, Wormholes and Superstrings.
\end{enumerate}

\end{document}